\section{Methodology}

\subsection{Problem Formulation}
We consider a Markov Decision Process (MDP) defined by $(\gS, \gA, R, P, \gamma)$, where $\gS$ and $\gA$ are the state and action spaces, $R(s, a)$ is the reward function, $P(s'|s, a)$ is the transition probability, and $\gamma \in [0, 1)$ is the discount factor. In distributional RL, we are interested in modeling the random variable $Z^\pi(s, a)$ representing the discounted sum of future rewards (the return) under policy $\pi$:
\begin{equation}
Z^\pi(s, a) \stackrel{D}{=} R(s, a) + \gamma Z^\pi(s', a'), \quad s' \sim P(\cdot|s, a), \ a' \sim \pi(\cdot|s')
\end{equation}
where $\stackrel{D}{=}$ denotes equality in distribution. The distributional Bellman operator $\gT^\pi$ is defined as $\gT^\pi Z(s, a) := R(s, a) + \gamma Z(s', a')$.

\subsection{Inductive Moment Matching Q-learning (\immq)}
Our proposed method, \immq, approximates the distribution $Z(s, a)$ using a set of $N$ unconstrained particles $\{\theta_i(s, a)\}_{i=1}^N$. A neural network $Q_\phi(s, a)$ outputs these $N$ particles. To update the network, we minimize the Maximum Mean Discrepancy (\mmd) between the predicted particles and the target particles derived from the Bellman update.

\para{The IMM Objective} Building on the \imm framework \cite{zhou2025inductive}, we define the loss function for a transition $(s, a, r, s')$ as:
\begin{equation}
\gL(\phi) = \text{MMD}^2(\{z_i\}_{i=1}^N, \{\hat{z}_j\}_{j=1}^N) + \alpha \gL_{\text{ind}}(z, \hat{z})
\end{equation}
where $z_i = Q_\phi(s, a)$ are the current predicted particles, and $\hat{z}_j = r + \gamma Q_{\bar{\phi}}(s', a')$ are the target particles generated by a target network $Q_{\bar{\phi}}$. The \mmd distance is computed using a sum of multiple RBF kernels with varying bandwidths $\sigma \in \{0.1, 1.0, 10.0\}$ to capture different scales of the distribution.

The inductive term $\gL_{\text{ind}}$ regularizes the mean of the predicted distribution toward the mean of the target distribution:
\begin{equation}
\gL_{\text{ind}} = \left\| \frac{1}{N} \sum_{i=1}^N z_i - \frac{1}{N} \sum_{j=1}^N \hat{z}_j \right\|^2_2
\end{equation}
The hyperparameter $\alpha$ controls the strength of this inductive bias. When $\alpha=0$, the method reduces to a standard particle-based \mmd update. For $\alpha > 0$, the network is encouraged to match the first moment (the mean) explicitly, allowing the \mmd term to focus on matching higher-order moments and the overall shape of the distribution.

\subsection{Implementation Details}
We implement \immq using a feed-forward neural network with two hidden layers of 256 units each. The output layer predicts $N=32$ particles. We use the Adam optimizer with a learning rate of $3 \times 10^{-4}$ and perform soft target updates with $\tau = 0.005$. Training is conducted in an offline setting using the `halfcheetah_medium.hdf5` dataset from \dfourrl \cite{fu2020d4rl}. We split the selected trajectories into 80\% training and 20\% validation.
